\documentclass[10pt]{article}
\usepackage[a4paper,margin=2.2cm]{geometry}
\usepackage{lmodern,microtype,amsmath,amssymb,booktabs,caption,enumitem,siunitx,xcolor,hyperref}
\hypersetup{colorlinks,linkcolor=black,citecolor=black,urlcolor=blue!50!black}
\captionsetup{font=small,labelfont=bf}
\setlist{itemsep=1pt,topsep=3pt}
\newcommand{\qualityWorlds}{270}
\newcommand{\optimizeWorstGap}{1.5e-12}
\newcommand{\greedyPlainMedianGap}{9.98}
\newcommand{\greedyPlainWorstGap}{89.59}
\newcommand{\threads}{16}
\newcommand{\originalTableDeviation}{4e-13}
\newcommand{\maxContenders}{4}
\newcommand{\bianchiErrorScenario}{0.005}
\newcommand{\bianchiErrorLarge}{0.018}
\newcommand{\hiddenError}{0.002}
\newcommand{\alohaError}{0.218}
\newcommand{\estimateLossError}{0.49}
\newcommand{\estimateCodecError}{0.10}
\newcommand{\isolatedCases}{48}
\newcommand{\isolatedLossError}{0.17}
\newcommand{\isolatedCodecError}{0.03}
\newcommand{\voiceExactCases}{80}
\newcommand{\voiceExactGap}{5e-13}
\newcommand{\hOneWorlds}{50}
\newcommand{\hOneWins}{49}
\newcommand{\hOneTies}{1}
\newcommand{\hOneLosses}{0}
\newcommand{\hOneMedian}{3.12}
\newcommand{\hOneP}{1.8e-15}
\newcommand{\hOneBWorlds}{50}
\newcommand{\hOneBWins}{50}
\newcommand{\hOneBTies}{0}
\newcommand{\hOneBLosses}{0}
\newcommand{\hOneBMedian}{136.40}
\newcommand{\hOneBP}{8.9e-16}
\newcommand{\hTwoWorlds}{50}
\newcommand{\hTwoWins}{50}
\newcommand{\hTwoTies}{0}
\newcommand{\hTwoLosses}{0}
\newcommand{\hTwoMedian}{27.19}
\newcommand{\hTwoP}{8.9e-16}
\newcommand{\hThreeKept}{200}
\newcommand{\hThreeTrials}{200}
\newcommand{\hTwoCodecCut}{26.1}

\newcommand{\code}[1]{\texttt{#1}}

\title{Channel assignment and voice QoS in synthetic 802.11 networks:\\ model, methods and verification}
\author{}
\date{\today}

\begin{document}
\maketitle

\begin{abstract}
\noindent This note describes a Julia model (\code{WiFiAssignment.jl}) that assigns channels to 802.11 access points so
as to maximise proportional fairness of client throughput, optionally together with the quality of voice calls carried
over the same network. It covers the radio and airtime model, the subset-table formulation that makes exact optimisation
possible for up to 18 access points, the heuristics, the voice queue and MAC model, the voice-aware objective, and how
each component is checked against an independent reference. Results are in the companion \code{results.pdf}.
\end{abstract}

\section{Network model}

\paragraph{Geometry and propagation.}
A world has $n \le 18$ access points (APs) placed uniformly in an $A \times A$ square, each serving $m$ clients placed at a
uniform distance of 5--15\,m and uniform angle from it (clipped to the square). Received power over distance $d$ is
\[
  P\ [\mathrm{dBm}] = p_0 - \alpha \log_{10}\max(d, 1) + \sigma Z, \qquad Z \sim \mathcal N(0,1),
\]
drawn independently per AP--client link and once per AP pair (symmetric). Defaults are $p_0 = -30$\,dBm, $\alpha = 35$
(path-loss exponent 3.5), $\sigma = 4$\,dB, noise $N_0 = -95$\,dBm, bandwidth $B = 20$\,MHz. Two APs \emph{sense} each
other when their mutual received power is at least the carrier-sense threshold ($-82$\,dBm); co-channel APs that do not
sense each other are \emph{hidden} from one another.

\paragraph{Client rates.}
An assignment $x$ gives each AP one of $k$ orthogonal channels. For client $c$ of AP $j$, with signal $S_c$, let
$\mathcal S_j(x)$ and $\mathcal H_j(x)$ be the co-channel APs that $j$ senses and does not sense, and $a_i \in [0,1]$ the
airtime of AP $i$. Hidden APs interfere in proportion to their airtime; sensed APs share airtime:
\begin{align*}
  \phi_c &= \min\!\Big(R_{\max},\; B \log_2\Big(1 + \frac{S_c}{N_0 + \sum_{h \in \mathcal H_j(x)} a_h P_{c,h}}\Big)\Big), \\
  r_c &= \min\!\Big(D,\; \max\!\Big(\frac{\eta\, \phi_c}{m\,\big(1 + \sum_{s \in \mathcal S_j(x)} a_s\big)},\; r_{\min}\Big)\Big),
\end{align*}
with PHY cap $R_{\max} = 144$\,Mbps, MAC efficiency $\eta = 0.72$, per-client demand $D = 40$\,Mbps and floor
$r_{\min} = 0.5$\,Mbps. The objective is proportional fairness \cite{kelly},
\[
  U(x) = \sum_c \log r_c \quad \text{(nats)}.
\]

\paragraph{Airtime as a fixed point.}
An AP's airtime is what its clients' demand needs at the rates they actually get:
\[
  a_j = \min\!\Big(1,\; \sum_{c \in j} \frac{D}{\eta \max(\phi_c(a), r_{\min})}\Big).
\]
Because $\phi_c$ falls as hidden APs' airtime rises, the right-hand side is monotone in $a$, so iterating from the
interference-free airtime converges to the least fixed point. When every AP already needs all its airtime (the default
demand with three or more clients per AP), the fixed point is $a \equiv 1$ and the iteration is skipped. Like every
iterative solver in the code, it raises an error rather than return an unconverged value.

\section{Optimisation}

\paragraph{Separability.}
Every term in $r_c$ depends only on the set $G$ of APs sharing AP $j$'s channel, so
$U(x) = \sum_{\text{channels}} f(G_{\text{ch}})$ with $f(G) = \sum_{j \in G} \sum_{c \in j} \log r_c(G)$. The $2^n$ values of
$f$ are tabulated once; any assignment then costs $k$ lookups and moving one AP costs four.

\paragraph{Building the table.}
For saturated worlds interference and blocked airtime add over APs, so each client needs subset-sum tables over the low
and high halves of the AP index (and each AP one table of blocked airtime); entries are formed row by row in parallel,
with up to 16 client rates multiplied before one logarithm. For unsaturated worlds each entry solves its own airtime
fixed point. Both paths are deterministic and give bit-identical results for any thread count.

\paragraph{Heuristics.}
\begin{itemize}
\item \emph{Graph colouring} (baseline): conflict graph at $-75$\,dBm, APs in decreasing degree, least-conflicting colour.
\item \emph{Greedy}: APs in decreasing order of the power they put on other APs' clients, each placed on the channel with
  the largest marginal table gain; optionally followed by \emph{local search} over single-AP moves and pairwise swaps until
  neither improves the objective.
\item \emph{Simulated annealing}: random start, random single-AP moves with Metropolis acceptance and geometric cooling to
  $T_0/100$. $T_0$ is calibrated so the mean downhill move from the start is accepted with probability $1/2$ (from 200
  sampled moves), and the best assignment is finished with local search.
\item \code{optimize\_assignment}: iterated local search. Eight restarts (the greedy order, then random orders), each
  followed by 200 kicks that reassign 2--4 random APs and re-run local search, keeping the result when it is no worse.
\end{itemize}

\paragraph{Exact search.}
Let $B_j(M)$ be the best split of AP set $M$ into at most $j$ channels. The channel holding the lowest AP of $M$ is some
$G \subseteq M$, so
\[
  B_1(M) = f(M), \qquad B_j(M) = \max_{G \ni \min M,\; G \subseteq M} \big[f(G) + B_{j-1}(M \setminus G)\big],
\]
and the optimum is $B_k(\{1,\dots,n\})$. AP 1 always opens the first channel, so the intermediate layers only need sets
without it, and the top layer only the full set: about $(k-2)\,3^{n-1}/2$ operations, parallel over sets. When $k \ge n$
every AP gets its own channel, which is optimal because adding an AP to a channel can only lower its members' rates.

\section{Voice model}

\paragraph{Traffic and queues.}
Each AP carries $K$ two-way calls, i.e.\ $2K$ flows that each send one packet per 20\,ms frame of
$b\cdot 20\,\text{ms} + 480$ bits at codec rate $b$ (64\,kbps nominal, 12\,kbps minimum), with a 150\,ms deadline.
An AP's raw capacity per frame is the sum of its clients' rates. Background traffic offers a fraction $\rho$ of that
capacity. Under FIFO, voice and background of a frame interleave, so each frame enters the queue as one mixed item and voice
loses in proportion to its share when the queue overflows; under priority scheduling only voice is queued. Work older than
the deadline is dropped from the head of the queue, and delay is measured at the time the packet's bits finish within the
frame.

\paragraph{MAC.}
AP $j$'s utilisation is $u_j = \min(1, \rho + \text{voice bits}/\text{raw capacity})$. A transmission collides with
probability
\[
  p_j = 1 - \big(1 - p^{\mathrm B}(1 + \textstyle\sum_{s \in \mathcal S_j} u_s)\big)
            \prod_{h \in \mathcal H_j} \big(1 - \sigma_{jh}\, q(u_h)\big),
\]
where $p^{\mathrm B}(n)$ is Bianchi's DCF fixed point \cite{bianchi} for $n$ contenders,
$p = 1 - (1-\tau)^{n-1}$, $\tau = 2/(1 + W + pW\sum_{i<m}(2p)^i)$ with $W = 16$ and $m = 6$ backoff stages (802.11 OFDM),
solved by a bracketed Newton iteration. A hidden AP busy a fraction $u$ of the time overlaps a frame with probability
$q(u) = u + (1-u)\big(1 - e^{-u/(1-u)}\big)$ (busy at its start, or starting during it, with frames of equal length and
exponential idle gaps), and $\sigma_{jh}$ is the share of $j$'s clients whose signal is within the 10\,dB capture
threshold of $h$'s. With $R = 2$ retries, service is capacity divided by $\sum_{i=0}^{R} p_j^i$ expected attempts, and a
fraction $p_j^{R+1}$ of served packets is lost.

\paragraph{Policies.}
\code{:fifo} and \code{:deadline} use the nominal codec; \code{:adaptive} (FIFO) and \code{:combined} (priority) move the codec
35\,\% of the way per frame towards the rate at which voice fits the airtime left to it, the service minus background under
FIFO and the whole service under priority, so the codec only shrinks when voice does not fit.

\paragraph{Voice-aware objective.}
An AP's voice outcome depends only on the APs sharing its channel, so it adds a second term to $f(G)$. For each $G$ the
steady state of the simulator (utilisation, collisions and codec rate) is solved as a joint fixed point, and each flow
contributes $\log(b/b_{\text{nominal}} \cdot \text{delivered fraction})$, so codec quality and loss both count and data
clients and voice flows are treated alike by proportional fairness. The same heuristics and exact search then apply
unchanged.

\section{Verification}

Each component is checked against a reference that shares no code with it. The counts below come from the
deterministic checks in \code{benchmarks/stress.jl} and \code{benchmarks/dense.jl}; the other figures are produced
by \code{paper/results.jl}.
\begin{itemize}
\item rates, scores and the airtime fixed point against an independent dB-domain implementation with its own Jacobi
  iteration: exhaustive search over 1--5 channels and up to 6 APs (5\,850 checks), every table entry of 40 worlds
  (8\,304), and 96 worlds of 1--18 APs (14\,112);
\item exact search and \code{optimize\_assignment} against an independent subset dynamic program for 2--5 channels
  (8\,376 checks for 3 channels, 9\,936 for 2, 4 and 5) and against each other at 18 APs for 3, 6, 12 and 18 channels;
\item the voice queue against the closed-form steady state of an isolated AP (\isolatedCases{} cases, loss within
  \isolatedLossError{} points), and packet accounting on 40 random worlds;
\item Bianchi's model against a slot-level DCF simulation (within \bianchiErrorScenario{} up to \maxContenders{} contenders,
  the most the voice scenarios reach) and the hidden overlap formula against simulation (within \hiddenError{});
\item the steady-state voice estimate against the simulator (loss within \estimateLossError{} points, codec within
  \estimateCodecError{} kbps), and the voice-aware optimizer against exhaustive search (\voiceExactCases{} cases).
\end{itemize}

\section{Changes made to the original code}

\begin{itemize}
\item \emph{Subset table}: rebuilt around per-client half-mask subset sums, one logarithm per group of client rates and
  parallel rows; it reproduces the original table to within \num{\originalTableDeviation}.
\item \emph{Heuristics}: local search after greedy, calibrated temperature and final polish for annealing, and the new
  iterated local search \code{optimize\_assignment}.
\item \emph{Exact search}: from three channels and 14 APs to any number of channels and 18 APs.
\item \emph{Airtime}: from the interference-free demand to the fixed point above.
\item \emph{Voice}: the adaptive codec now accounts for background traffic and only shrinks when voice does not fit;
  background work expires instead of accumulating; FIFO mixes voice and background in proportion instead of placing voice
  in the middle of each frame; delay is measured within the frame instead of at frame boundaries; the collision model
  $1 - e^{-0.12\,|\mathcal H_j|\,u_j}$ (capped at 0.65) is replaced by Bianchi contention plus capture-weighted
  hidden-terminal overlap.
\item \emph{Voice-aware assignment} and loud failure of every iterative solver.
\end{itemize}

\section{Limitations}
Worlds, traffic and shadowing are synthetic and static, and channels are orthogonal. The voice model is an aggregate work
queue rather than a packet-level 802.11 simulation. Bianchi assumes saturated contenders, so partly busy APs enter as
fractional contenders, and each AP's own clients count as a single contender, which likely understates collisions in
call-heavy cells. The exact methods are limited to 18 APs by the $2^n$ table. Conclusions hold for this model; external
validation needs ns-3, packet traces or hardware.

\begin{thebibliography}{9}
\bibitem{bianchi} G.~Bianchi, ``Performance analysis of the IEEE 802.11 distributed coordination function,''
  \emph{IEEE Journal on Selected Areas in Communications}, vol.~18, no.~3, pp.~535--547, 2000.
\bibitem{kelly} F.~Kelly, ``Charging and rate control for elastic traffic,''
  \emph{European Transactions on Telecommunications}, vol.~8, no.~1, pp.~33--37, 1997.
\end{thebibliography}

\end{document}
